\section{Introduction}
The modernization of the electrical power infrastructure into the ``Smart Grid'' represents a paradigm shift from centralized, one-way transmission to a decentralized, bi-directional network driven by data. To manage the increasing complexity of renewable energy integration and dynamic load balancing, operators are rapidly adopting Artificial Intelligence (AI) and Machine Learning (ML) techniques. These models are deployed for critical tasks such as load forecasting, fault detection, and real-time pricing optimization \cite{sanchez2024attacking}.

Despite their operational benefits, the reliance on deep learning models expands the grid's attack surface. Existing research in adversarial machine learning has predominantly focused on \textit{integrity} attacks, where an adversary perturbs input data (e.g., adversarial examples) to cause incorrect model predictions. However, in safety-critical systems like the Smart Grid, \textit{availability} is equally paramount. Decisions regarding frequency regulation and fault isolation must be made within milliseconds to prevent cascading failures.

This report addresses a specific threat to availability known as the ``Sponge Attack,'' formally introduced by Shumailov et al. \cite{shumailov2021sponge}. Unlike standard adversarial attacks that seek to fool a model, sponge attacks aim to ``soak up'' computational resources. By crafting specific inputs (Sponge Examples) \cite{cina2022sponge} or poisoning model weights during training (Sponge Poisoning) \cite{cina2022sponge}, an attacker can trigger the worst-case computational performance of a neural network. This results in significant spikes in inference latency and energy consumption.

The implications for the Smart Grid are severe. A sponge attack targeting a central load forecasting server could introduce latency that desynchronizes grid control loops. Furthermore, attacks on edge devices—such as Advanced Metering Infrastructure (AMI) or IoT sensors—can rapidly drain batteries, rendering the monitoring infrastructure blind.

To support reproducible research in Smart Grid security, we release our code, pre-trained ACT-LSTM models, and the XAI-guided attack pipeline at a public repository\footnote{\url{https://github.com/bbpbbp09/Sponge-Attack-Seminar}}.

This paper is structured to address three primary objectives:
\begin{itemize}
    \item \textbf{Literature Review:} A meta-analysis of AI-based smart grid systems and the mechanics of sponge attacks (Section \ref{sec:related}).
    \item \textbf{System Model:} A formal definition of the threat landscape and target architectures (Section \ref{sec:threat_model}).
    \item \textbf{Experimentation:} A rigorous comparison of Black-Box vs. White-Box attacks, including a novel bit-flip oracle that targets energy consumption via switching activity, and an evaluation of defense mechanisms (Section \ref{sec:experiments} and \ref{sec:results}).
\end{itemize}
